% Zone „Das kennst du schon“ – aus bank/rationale-zahlen/zone.jsonl (zusammenbau v0.1)
\einheitenkopf[zone][Kennst du schon]{Das kennst du schon}
\setcounter{aufgabe}{0}
%% TODO Grundvorstellungs-Aufgabe als erste Hauptnummer der Zone (2.8) fehlt in der Bank

%% TODO Titel: Ich-kann-Satz fehlt in der Bank (Platzhalter: Kettenname) – Nr. 1
%% TODO Anweisung: ein Satz über der ersten Teilaufgabe fehlt in der Bank – Nr. 1
\begin{aufgabe}{Natürliche Zahlen addieren, subtrahieren, multiplizieren, dividieren im Kopf (Einmaleins)}
%% TODO Darstellung für schwach fehlt in der Bank (rationale-zahlen-zone-f1-v1; Abschnitt „Für schwache Schüler“)
\swz{Berechne: $7 \cdot 8$}{}
%% TODO Darstellung für schwach fehlt in der Bank (rationale-zahlen-zone-f1-v3; Abschnitt „Für schwache Schüler“)
\swz{Berechne: $120 : 8$}{}
\end{aufgabe}

%% TODO Titel: Ich-kann-Satz fehlt in der Bank (Platzhalter: Kettenname) – Nr. 2
%% TODO Anweisung: ein Satz über der ersten Teilaufgabe fehlt in der Bank – Nr. 2
\begin{aufgabe}{Zahlenstrahl}
\swz{Welche Zahl gehört zu A? Lies ab.}{\zahlenstrahl[xmin=0,xmax=10,xstep=1]{7/A}}
\swz{Welche Zahl gehört zu C? Lies ab.}{\zahlenstrahl[xmin=0,xmax=2,xstep=0.1,karo=0.5]{1.3/C}}
\end{aufgabe}

%% TODO Titel: Ich-kann-Satz fehlt in der Bank (Platzhalter: Kettenname) – Nr. 3
%% TODO Anweisung: ein Satz über der ersten Teilaufgabe fehlt in der Bank – Nr. 3
\begin{aufgabe}{Dezimalzahlen und Brüche addieren, subtrahieren, multiplizieren (Kommazahlen addieren, Bruch mal ganze Zahl)}
%% TODO Darstellung für schwach fehlt in der Bank (rationale-zahlen-zone-f3-v1; Abschnitt „Für schwache Schüler“)
\swz{Berechne: $2{,}5 + 1{,}3$}{}
%% TODO Darstellung für schwach fehlt in der Bank (rationale-zahlen-zone-f3-v3; Abschnitt „Für schwache Schüler“)
\swz{Berechne: $4{,}6 - 1{,}85$}{}
\end{aufgabe}

%% TODO Titel: Ich-kann-Satz fehlt in der Bank (Platzhalter: Kettenname) – Nr. 4
%% TODO Anweisung: ein Satz über der ersten Teilaufgabe fehlt in der Bank – Nr. 4
\begin{aufgabe}{Punkt vor Strich und Klammern mit natürlichen Zahlen}
%% TODO Darstellung für schwach fehlt in der Bank (rationale-zahlen-zone-f4-v1; Abschnitt „Für schwache Schüler“)
\swz{Berechne: $3 + 4 \cdot 5$}{}
%% TODO Darstellung für schwach fehlt in der Bank (rationale-zahlen-zone-f4-v3; Abschnitt „Für schwache Schüler“)
\swz{Berechne: $0{,}5 + 4 \cdot 2{,}5$}{}
\end{aufgabe}

%% TODO Titel: Ich-kann-Satz fehlt in der Bank (Platzhalter: Kettenname) – Nr. 5
%% TODO Anweisung: ein Satz über der ersten Teilaufgabe fehlt in der Bank – Nr. 5
\begin{aufgabe}{Wert eines Terms mit Platzhalter berechnen (Vorzahl mal x plus Zahl, für einen gegebenen Einsetzwert)}
%% TODO Darstellung für schwach fehlt in der Bank (rationale-zahlen-zone-f5-v1; Abschnitt „Für schwache Schüler“)
\swz{$3 \cdot x + 2$ für $x = 4$ – welcher Wert?}{}
%% TODO Darstellung für schwach fehlt in der Bank (rationale-zahlen-zone-f5-v3; Abschnitt „Für schwache Schüler“)
\swz{$2 \cdot x + 1$ für $x = 2{,}5$ – welcher Wert?}{}
\end{aufgabe}

%% TODO Titel: Ich-kann-Satz fehlt in der Bank (Platzhalter: Kettenname) – Nr. 6
%% TODO Anweisung: ein Satz über der ersten Teilaufgabe fehlt in der Bank – Nr. 6
\begin{aufgabe}{Punkt vor Strich und Klammern mit natürlichen Zahlen – Fehler finden}
\swz[3]{Mia rechnet: $27 - 9 : 3 = 18 : 3 = 6$. Finde den Fehler.}{}
\end{aufgabe}

%% TODO Titel: Ich-kann-Satz fehlt in der Bank (Platzhalter: Kettenname) – Nr. 7
%% TODO Anweisung: ein Satz über der ersten Teilaufgabe fehlt in der Bank – Nr. 7
\begin{aufgabe}{Punkt vor Strich und Klammern mit natürlichen Zahlen – selbst rechnen}
%% TODO Darstellung für schwach fehlt in der Bank (rationale-zahlen-zone-f4-v6; Abschnitt „Für schwache Schüler“)
\swz{Berechne: $40 - 18 : 6$}{}
\end{aufgabe}
